% ===================================================================
%  ТАБЛИЦА № 3 — Детство и Юность.
%  Traduction 2026 d'apres texto/io, controlee sur texto/fr, texto/en
%  et texto/es. Consigne : voir texto/ru/10-tablica-01.tex.
%
%  Deux scenes, et l'ido subdivise beaucoup plus que le francais :
%  huit des intertitres n'ont pas de vis-a-vis francais. Le russe suit
%  l'ido, comme les autres traductions.
% ===================================================================

%%K t03-apar-1 apar
\VUsaut{3.0mm}
\VUtitre{15.0pt}{60}{ТАБЛИЦА № 3}
\VUsaut{1.6mm}
\VUtitre{11.4pt}{0}{Детство и Юность.}
\VUsaut{3.4mm}

%%K t03-apar-2 apar
(Таблицы 3 и 4 составлены так, чтобы дать в четырёх \VUgras{семейных
сценах} основной запас слов о \VUgras{погоде}, \VUgras{возрасте},
\VUgras{семье}, \VUgras{одежде} и \VUgras{здоровье}.)
\VUsaut{4.0mm}

%%K t03-tit-1 sub
\VUcentre{12.0pt}{0}{\textit{Первая сцена.}}
\VUsaut{3.0mm}

%%K t03-tit-2 sub
\VUpk{10.2pt}{40}{Крестины в деревне.}
\VUsaut{2.4mm}

%%K t03-tit-3 sub
\VUpk{10.2pt}{40}{Весна.}
\VUsaut{3.0mm}

%%K t03-c1-01-1 p
1. --- Мы в \VUgras{весне}, в месяце \VUgras{марте}, \VUgras{апреле} или
\VUgras{мае}, в один из дней \VUgras{недели} --- в \VUgras{понедельник},
\VUgras{вторник} или \VUgras{среду}. \VUgras{Десять часов} утра.

%%K t03-c1-02-1 p
2. --- \VUgras{Погода} прекрасная, \VUgras{воздух} чистый. Светит солнце.
Несколько маленьких \VUgras{облаков} \textsuperscript{(2)} поднимаются в синем
\VUgras{небе} \textsuperscript{(1)}.

%%K t03-c1-03-1 p
3. --- На \VUgras{деревьях} почти нет \VUgras{листьев}. Стройные
\VUgras{тополя} \textsuperscript{(3)} и высокий \VUgras{платан} \textsuperscript{(17)}
покрыты пока только почками.

%%K t03-c1-04-1 p
4. --- \VUgras{Ласточки} \textsuperscript{(4)} вернулись и с радостными криками
кружат вокруг \VUgras{шпиля} \textsuperscript{(7)} \VUgras{колокольни}
\textsuperscript{(5)}, венчающей деревенскую \VUgras{церковь} \textsuperscript{(6)}.

%%K t03-tit-4 sub
\VUsaut{3.4mm}
\VUpk{10.2pt}{40}{Церковь.}
\VUsaut{3.0mm}

%%K t03-c1-05-1 p
5. --- Очень проста эта церковь с большим \VUgras{порталом} \textsuperscript{(13)}
и большими \VUgras{часами} \textsuperscript{(8)}. Малая \VUgras{стрелка}
\textsuperscript{(9)} стоит на цифре X: она показывает \VUgras{часы}.
\VUgras{Большая} \textsuperscript{(10)} стоит на XII (полдень); она показывает
\VUgras{минуты}.

%%K t03-c1-06-1 p
6. --- На колокольне весело звонят \VUgras{колокола} \textsuperscript{(11)}. На
верху крыши стоит небольшой каменный \VUgras{крест} \textsuperscript{(12)}.

%%K t03-c1-07-1 p
7. --- У церкви не только большой \VUgras{портал} \textsuperscript{(13)}, есть и
маленькая боковая дверь. Через эту дверь \VUgras{священник}
\textsuperscript{(15)} в \VUgras{сутане} выходит из \VUgras{ризницы}
\textsuperscript{(14)} и направляется в \VUgras{дом священника} \textsuperscript{(19)}.

%%K t03-c1-08-1 p
8. --- Рядом с ним --- \VUgras{молочница} \textsuperscript{(24)} со своим
\VUgras{ослом} \textsuperscript{(23)}. Она возвращается из города, куда возила
детям своё доброе молоко.

%%K t03-tit-5 sub
\VUsaut{4.0mm}
\VUpk{10.2pt}{40}{Фруктовый сад.}
\VUsaut{3.4mm}

%%K t03-c1-09-1 p
9. --- Господин Серый (осёл) очень доволен этой утренней прогулкой. Он с
наслаждением вдыхает воздух, надушенный \VUgras{цветом} \textsuperscript{(20)},
который покрывает \VUgras{фруктовые деревья} \textsuperscript{(18)} соседнего
\VUgras{сада}. Все в розовом и белом эти \VUgras{персиковые деревья}
\textsuperscript{(18)} и \VUgras{абрикосовые деревья} \textsuperscript{(19)}, и они
обещают хороший урожай.

%%K t03-c1-10-1 p
10. --- А тем временем маленькие \VUgras{пчёлы} \textsuperscript{(22)} из
\VUgras{улья} \textsuperscript{(21)} летят туда собирать свой сладкий
\VUgras{мёд}, жужжа за работой.

%%K t03-c1-11-1 p
11. --- Но что случилось? Вот бегут деревенские дети и разгоняют
испуганных \VUgras{гусей} \textsuperscript{(25)}. Это шествие выходит из церкви
после \VUgras{крестин} сына нотариуса.

%%K t03-c1-12-1 p
12. --- Маленький Яков в своей \VUgras{колыбели} \textsuperscript{(26)}
просыпается от весёлого звона колоколов, а его сестра Магдалина, которой
тоже хочется посмотреть на шествие, просит мать причесать и одеть её на
пороге.

%%K t03-c1-13-1 p
13. --- Мать соглашается. Она надевает свой \VUgras{платок}
\textsuperscript{(28)}, свою \VUgras{кофту} \textsuperscript{(29)} и свой
\VUgras{передник} \textsuperscript{(30)} и садится на низкий стул перед дверью.
На Магдалине только \VUgras{юбка} \textsuperscript{(31)} и \VUgras{лифчик}
\textsuperscript{(32)}.

%%K t03-c1-14-1 p
14. --- Как она счастлива! Она забывает свои любимые цветы, свои
\VUgras{гвоздики} \textsuperscript{(34)}, на которые садятся \VUgras{бабочки}
\textsuperscript{(35)}, свои \VUgras{тюльпаны} \textsuperscript{(36)}, свои
\VUgras{ландыши} \textsuperscript{(37)} в \VUgras{горшках} \textsuperscript{(38)} на
\VUgras{стене} \textsuperscript{(33)} и свои \VUgras{розы} \textsuperscript{(39)},
которые образуют такую красивую изгородь. Она заглядывается на богатые
наряды дам в шествии.

%%K t03-c1-15-1 p
15. --- Деревенские мальчики мало смотрят на наряды. Они бросаются
вперёд и толкают друг друга, чтобы получить \VUgras{конфеты}
\textsuperscript{(55)} или \VUgras{монеты} \textsuperscript{(52)}. Один из них плачет
\textsuperscript{(40)}.

%%K t03-c1-16-1 p
16. --- Другой наступил на лапу \VUgras{собаке} \textsuperscript{(42)}, и та с
визгом убегает и обращает в бегство \VUgras{курицу} \textsuperscript{(43)} с её
\VUgras{цыплятами} \textsuperscript{(44)}.

%%K t03-tit-6 sub
\VUsaut{5.0mm}
\VUpk{10.2pt}{40}{Крестильное шествие.}
\VUsaut{4.0mm}

%%K t03-c1-17-1 p
17. --- Во главе шествия идёт \VUgras{кормилица} \textsuperscript{(45)}. На ней
красивый \VUgras{чепец} \textsuperscript{(46)} с длинными лентами. Она несёт
\VUgras{новорождённого} \textsuperscript{(47)}, всего в \VUgras{кружевах}
\textsuperscript{(48)}.

%%K t03-c1-18-1 p
18. --- За кормилицей идут \VUgras{крёстный} \textsuperscript{(49)} и
\VUgras{крёстная} \textsuperscript{(53)}. Они раздают конфеты и монеты. У
крёстного \VUgras{перчатки} \textsuperscript{(51)} и \VUgras{цилиндр}
\textsuperscript{(50)}. Крёстная держит в руке красивый \VUgras{букет}
\textsuperscript{(54)}.

%%K t03-c1-19-1 p
19. --- За ними видны \VUgras{дядя} \textsuperscript{(56)} и \VUgras{тётя}
\textsuperscript{(58)} ребёнка. На тёте изящная \VUgras{шляпа} \textsuperscript{(59)},
на дяде чёрный \VUgras{сюртук} \textsuperscript{(57)} и белый жилет. Дальше идут
\VUgras{двоюродный брат} \textsuperscript{(60)} и \VUgras{двоюродная сестра}
\textsuperscript{(61)}. Последняя ведёт за руку \VUgras{сестру} \textsuperscript{(62)}
новорождённого, маленькую Берту.

%%K t03-c1-20-1 p
20. --- Другой \VUgras{брат} \textsuperscript{(63)} Берты идёт позади. Он подаёт
руку \VUgras{родственнице} \textsuperscript{(64)}. За ними идут \VUgras{друзья}
\textsuperscript{(65)} семьи.

%%K t03-tit-7 sub
\VUsaut{4.6mm}
\VUpk{10.2pt}{40}{Зеваки.}
\VUsaut{3.4mm}

%%K t03-c1-21-1 p
21. --- Возле шествия --- \VUgras{нищий} \textsuperscript{(66)}, опирающийся на
свой \VUgras{костыль} \textsuperscript{(67)}. Он просит милостыни.

%%K t03-c1-22-1 p
22. --- По другую сторону --- очень \VUgras{вежливый} \textsuperscript{(68)}
мальчик. Он приподнимает шапку и говорит \VUgras{«здравствуйте»} тем, кого
знает.

%%K t03-c1-23-1 p
23. --- Деревенские вышли из домов посмотреть на крестильное шествие.
Мы видим \VUgras{крестьянина} \textsuperscript{(69)} в \VUgras{бумажном колпаке},
а рядом с ним \VUgras{крестьянку} \textsuperscript{(70)}. Затем \VUgras{старуху}
\textsuperscript{(71)}, опирающуюся на свою \VUgras{палку} \textsuperscript{(72)}.
Наконец \VUgras{мельника} \textsuperscript{(73)} в широкой \VUgras{войлочной
шляпе} \textsuperscript{(74)} и молодую крестьянку в \VUgras{деревянных башмаках}
\textsuperscript{(75)}, которая учит ходить своего малыша.

%%K t03-c1-24-1 p
24. --- Картина очень оживлённая.
\VUsaut{4.0mm}

%%K t03-tit-8 sub
\VUcentre{12.0pt}{0}{\textit{Вторая сцена.}}
\VUsaut{3.0mm}

%%K t03-tit-9 sub
\VUpk{10.2pt}{40}{Народный праздник.}
\VUsaut{2.4mm}

%%K t03-tit-10 sub
\VUpk{10.2pt}{40}{Водные состязания и гонки.}
\VUsaut{3.0mm}

%%K t03-c2-01-1 p
1. --- Пришло \VUgras{лето}. Палящее солнце месяца \VUgras{июня} (июля или
\VUgras{августа}) манит молодёжь купаться и кататься на лодках.

%%K t03-c2-02-1 p
2. --- На правом берегу реки гребцы раздеваются, чтобы принять участие в
водных состязаниях и в гонках.

%%K t03-c2-03-1 p
3. --- Один из них снимает \VUgras{ботинки} \textsuperscript{(1)}; он
расшнуровывает (расстёгивает) их. Двое его товарищей уже готовы. На них
теперь только \VUgras{фуфайка} \textsuperscript{(2)} и \VUgras{плавки}
\textsuperscript{(3)}. Они болтают, поджидая товарищей. Они говорят о ярмарке
и о празднике на другом берегу.

%%K t03-c2-04-1 p
4. --- На земле рядом с ними лежит \VUgras{одежда} их товарищей:
\VUgras{пара} \VUgras{сапог} \textsuperscript{(4)}, \VUgras{туфли}
\textsuperscript{(5)}, \VUgras{носки} \textsuperscript{(6)}, \VUgras{войлочная}
\textsuperscript{(7)} и \VUgras{соломенная} \textsuperscript{(8)} шляпы и
\VUgras{галстук} \textsuperscript{(9)}.

%%K t03-c2-05-1 p
5. --- Один из юношей аккуратно сложил свой \VUgras{пиджак}
\textsuperscript{(10)}. Он кладёт его на полотенце, в которое его сосед сейчас
завернёт оба их \VUgras{костюма}.

%%K t03-c2-06-1 p
6. --- Шестой мальчик запаздывает. На голове у него ещё \VUgras{кепка}
\textsuperscript{(11)}. Он не снял ни \VUgras{рубашки} \textsuperscript{(12)}, ни
\VUgras{воротничка} \textsuperscript{(13)}, ни \VUgras{манжет} \textsuperscript{(14)}.
Он держит в руке свой \VUgras{жилет} \textsuperscript{(17)} и всё ещё в
\VUgras{брюках} \textsuperscript{(16)} и \VUgras{подтяжках} \textsuperscript{(15)}. К
следующему заезду он наверняка не поспеет.

%%K t03-c2-07-1 p
7. --- Возле юношей тропинка, обросшая \VUgras{чертополохом}
\textsuperscript{(18)}, спускается к воде, где старый Яков готовит среди
\VUgras{камышей} \textsuperscript{(19)} \VUgras{фейерверк} \textsuperscript{(20)},
который пустят вечером.

%%K t03-tit-11 sub
\VUsaut{4.4mm}
\VUpk{10.2pt}{40}{Гонка.}
\VUsaut{3.4mm}

%%K t03-c2-08-1 p
8. --- На реке несколько дам, укрывшись под зонтиками, катаются в
\VUgras{лодке} \textsuperscript{(21)}. Они следят за гонкой, а она очень
занимательна. В каждой \VUgras{гичке} \textsuperscript{(22)} четверо
\VUgras{гребцов} \textsuperscript{(24)} гребут изо всех сил, а \VUgras{рулевой}
\textsuperscript{(26)} держит \VUgras{руль} \textsuperscript{(25)} и правит хрупким
судёнышком. \VUgras{Вёсла} \textsuperscript{(23)} мерно бьют по воде, и лёгкие
лодки несутся с головокружительной быстротой.

%%K t03-c2-09-1 p
9. --- Несколько юношей состязаются возле быка моста за приз
\VUgras{мачты с призом} \textsuperscript{(28)}. Один из них осторожно
подвигается по упругому и скользкому шесту, чтобы достать
\VUgras{флажок} \textsuperscript{(29)}, укреплённый на конце. Его незадачливый
\VUgras{соперник} \textsuperscript{(30)} упал в воду. Он плывёт обратно к берегу.

%%K t03-c2-10-1 p
10. --- Ребята постарше \VUgras{бьются на воде} \textsuperscript{(32)} возле
\VUgras{набережной} \textsuperscript{(33)}. Большая \VUgras{толпа}
\textsuperscript{(31)} смотрит на все эти игры, облокотившись на \VUgras{перила}
\VUgras{моста} \textsuperscript{(27)} и теснясь вдоль \VUgras{берега}. Иные
зеваки, впрочем, оборачиваются, чтобы поглядеть на игру со
\VUgras{столбом} \textsuperscript{(45)} или полюбоваться \VUgras{шествием мэрии},
идущим на \VUgras{раздачу наград}.

%%K t03-c2-11-1 p
11. --- Сперва бегут несколько \VUgras{мальчишек} \textsuperscript{(35)} впереди
\VUgras{музыкантов} \textsuperscript{(36)}; затем идут \VUgras{мэр}
\textsuperscript{(37)}, его помощники и \VUgras{городской совет} городка.
Последними шагают \VUgras{пожарные} \textsuperscript{(38)}.

%%K t03-tit-12 sub
\VUsaut{5.0mm}
\VUpk{10.2pt}{40}{Ярмарка.}
\VUsaut{3.6mm}

%%K t03-c2-12-1 p
12. --- На \VUgras{ярмарочной площади} \textsuperscript{(39)} большая толпа.
Люди приходят посмотреть на разные \VUgras{балаганы}: на
\VUgras{деревянных лошадок} \textsuperscript{(40)}, на \VUgras{цирк}
\textsuperscript{(41)}, где представление уже началось; на \VUgras{народный бал}
\textsuperscript{(42)}, где молодые люди и девушки городка предаются
удовольствию \VUgras{танцевать}.

%%K t03-c2-13-1 p
13. --- В глубине видна \VUgras{арена} \textsuperscript{(44)}, где отважный
испанский \VUgras{матадор} бьётся с разъярённым \VUgras{быком}
\textsuperscript{(45)}; затем \VUgras{русские горы} \textsuperscript{(49)} и
\VUgras{тир} \textsuperscript{(46)}, где упражняются в обращении с
\VUgras{огнестрельным оружием} и в стрельбе по \VUgras{мишени}.

%%K t03-c2-14-1 p
14. --- Рядом --- \VUgras{балаганный театр} \textsuperscript{(47)}, где играют
\VUgras{комедию} и \VUgras{драму}. Совсем близко стоит балаган
\VUgras{великана} \textsuperscript{(48)} и \VUgras{карлика}. \VUgras{Рост}
первого --- два метра пятнадцать; второй едва с ребёнка.

%%K t03-c2-15-1 p
15. --- Наконец, вот большой \VUgras{зверинец} \textsuperscript{(50)} с его
\VUgras{дикими зверями}, с \VUgras{тигром} \textsuperscript{(51)} и его мощными
\VUgras{когтями}, со \VUgras{львом} \textsuperscript{(52)} и его густой
\VUgras{гривой}, с двумя \VUgras{жирафами} \textsuperscript{(53)} и
\VUgras{слоном} \textsuperscript{(54)}, который так ловко управляется со своим
\VUgras{хоботом}.

%%K t03-c2-16-1 p
16. --- \VUgras{Укротитель} \textsuperscript{(55)}, который обучает этих зверей,
стоит у двери балагана.
\VUsaut{6.0mm}
\VUfilet{22mm}
